% GENERATED FILE. Edit canonical lessons, then run npm run generate:books.
% canonical-source-hash: fnv1a64:e5fcb6a93e454059
% canonical-lessons: HI-C357-mulya, HI-C357-iman, HI-C357-izzat, HI-C357-astha, HI-C357-apnana

\chapter{What We Value}
\label{ch:hi-what-we-value}
\hlchaptermodality{357}

\begin{chapteropening}
\textbf{By the end of this chapter:} \emph{I can say what a family or a village values.}

Complete the last lesson of chapter 357: I can say what a family or a village values.
\end{chapteropening}

\section[\texorpdfstring{\hi{मूल्य} (mūlya)}{मूल्य (mūlya)}]{\hi{मूल्य} (mūlya) — a value (moral); worth}

\label{lesson:HI-C357-mulya}



\begin{quote}
Before the new one: say the Hindi for fairness, then the Hindi for to provoke.
\end{quote}

\subsection*{You'll want to know: \hi{मूल्य}}
\textbf{\hi{मूल्य}} — \emph{mūlya} — \textquotedblleft{}a value (moral); worth\textquotedblright{}.

\begin{quote}
\hi{सच} \hi{बोलना} \hi{हमारे} \hi{परिवार} \hi{का} \hi{सबसे} \hi{बड़ा} \hi{मूल्य} \hi{है।}
\end{quote}

\begin{quote}
\emph{sac bolnā hamāre parivār kā sabse baṛā mūlya hai.}
\end{quote}

\begin{quote}
Telling the truth is our family's greatest value.
\end{quote}

Masculine.

\begin{cousinweb}
\textbf{\hi{मूल्य}} (\emph{mūlya}) is a tatsama, borrowed from Sanskrit \emph{mūlya}, \textquotedblleft{}price, worth\textquotedblright{} (Monier-Williams). The moral sense, a value one lives by, is not among the senses Monier-Williams gives.
\end{cousinweb}

\begin{grammarlens}[title={What you've built}]
One more word for giving reasons and opinions.
\end{grammarlens}

\subsection*{Your turn}
\begin{itemize}
\raggedright
  \item \emph{Say it:} \emph{mūlya}
  \item \emph{Say it:} \emph{mūlya}, once more
  \item \emph{Say it:} \emph{sac bolnā hamāre parivār kā sabse baṛā mūlya hai.} — Telling the truth is our family's greatest value.
  \item \emph{Recall:} say the Hindi for fairness, then the Hindi for to provoke, then say \emph{mūlya} again
\end{itemize}

\begin{tcolorbox}[breakable,colback=teal!4,colframe=teal!35!black,title={Before you move on}]
What does \hi{मूल्य} mean? (\textquotedblleft{}A value (moral); worth\textquotedblright{}.) Say it once more.
\end{tcolorbox}
\section[\texorpdfstring{\hi{ईमान} (īmān)}{ईमान (īmān)}]{\hi{ईमान} (īmān) — integrity, good faith}

\label{lesson:HI-C357-iman}



\begin{quote}
Before the new one: say the Hindi for to provoke, then the Hindi for a value.
\end{quote}

\subsection*{You'll want to know: \hi{ईमान}}
\textbf{\hi{ईमान}} — \emph{īmān} — \textquotedblleft{}integrity, good faith\textquotedblright{}.

\begin{quote}
\hi{उसने} \hi{पैसे} \hi{के} \hi{लिए} \hi{अपना} \hi{ईमान} \hi{नहीं} \hi{बेचा।}
\end{quote}

\begin{quote}
\emph{usne paise ke liye apnā īmān nahīṃ becā.}
\end{quote}

\begin{quote}
He didn't sell his integrity for money.
\end{quote}

Masculine.

\begin{cousinweb}
\textbf{\hi{ईमान}} (\emph{īmān}) is usually traced to Arabic \emph{īmān}, \textquotedblleft{}faith\textquotedblright{}, which reached Hindi through Persian.
\end{cousinweb}

\begin{grammarlens}[title={What you've built}]
One more word for giving reasons and opinions.
\end{grammarlens}

\subsection*{Your turn}
\begin{itemize}
\raggedright
  \item \emph{Say it:} \emph{īmān}
  \item \emph{Say it:} \emph{īmān}, once more
  \item \emph{Say it:} \emph{usne paise ke liye apnā īmān nahīṃ becā.} — He didn't sell his integrity for money.
  \item \emph{Recall:} say the Hindi for to provoke, then the Hindi for a value, then say \emph{īmān} again
\end{itemize}

\begin{tcolorbox}[breakable,colback=teal!4,colframe=teal!35!black,title={Before you move on}]
What does \hi{ईमान} mean? (\textquotedblleft{}Integrity, good faith\textquotedblright{}.) Say it once more.
\end{tcolorbox}
\section[\texorpdfstring{\hi{इज़्ज़त} (izzat)}{इज़्ज़त (izzat)}]{\hi{इज़्ज़त} (izzat) — honour, standing}

\label{lesson:HI-C357-izzat}



\begin{quote}
Before the new one: say the Hindi for a value, then the Hindi for integrity.
\end{quote}

\subsection*{You'll want to know: \hi{इज़्ज़त}}
\textbf{\hi{इज़्ज़त}} — \emph{izzat} — \textquotedblleft{}honour, standing\textquotedblright{}.

\begin{quote}
\hi{गाँव} \hi{में} \hi{सब} \hi{उसकी} \hi{इज़्ज़त} \hi{करते} \hi{हैं।}
\end{quote}

\begin{quote}
\emph{gā̃v meṁ sab uskī izzat karte haiṁ.}
\end{quote}

\begin{quote}
Everyone in the village respects him.
\end{quote}

Feminine.

\begin{cousinweb}
\textbf{\hi{इज़्ज़त}} (\emph{izzat}) is usually traced to Arabic \emph{ʿizza}, \textquotedblleft{}might, honour\textquotedblright{}, which reached Hindi through Persian.
\end{cousinweb}

\begin{grammarlens}[title={What you've built}]
One more word for giving reasons and opinions.
\end{grammarlens}

\subsection*{Your turn}
\begin{itemize}
\raggedright
  \item \emph{Say it:} \emph{izzat}
  \item \emph{Say it:} \emph{izzat}, once more
  \item \emph{Say it:} \emph{gā̃v meṁ sab uskī izzat karte haiṁ.} — Everyone in the village respects him.
  \item \emph{Recall:} say the Hindi for a value, then the Hindi for integrity, then say \emph{izzat} again
\end{itemize}

\begin{tcolorbox}[breakable,colback=teal!4,colframe=teal!35!black,title={Before you move on}]
What does \hi{इज़्ज़त} mean? (\textquotedblleft{}Honour, standing\textquotedblright{}.) Say it once more.
\end{tcolorbox}
\section[\texorpdfstring{\hi{आस्था} (āsthā)}{आस्था (āsthā)}]{\hi{आस्था} (āsthā) — belief, steady faith (in something)}

\label{lesson:HI-C357-astha}



\begin{quote}
Before the new one: say the Hindi for integrity, then the Hindi for honour.
\end{quote}

\subsection*{You'll want to know: \hi{आस्था}}
\textbf{\hi{आस्था}} — \emph{āsthā} — \textquotedblleft{}belief, steady faith (in something)\textquotedblright{}.

\begin{quote}
\hi{मेरी} \hi{दादी} \hi{की} \hi{पूजा} \hi{में} \hi{गहरी} \hi{आस्था} \hi{है।}
\end{quote}

\begin{quote}
\emph{merī dādī kī pūjā meṁ gahrī āsthā hai.}
\end{quote}

\begin{quote}
My grandmother has deep faith in worship.
\end{quote}

Feminine.

\begin{cousinweb}
\textbf{\hi{आस्था}} (\emph{āsthā}) is a tatsama, borrowed from Sanskrit \emph{āsthā}, \textquotedblleft{}regard, care for\textquotedblright{} (Monier-Williams). Hindi uses it for steady faith in something.
\end{cousinweb}

\begin{grammarlens}[title={What you've built}]
One more word for giving reasons and opinions.
\end{grammarlens}

\subsection*{Your turn}
\begin{itemize}
\raggedright
  \item \emph{Say it:} \emph{āsthā}
  \item \emph{Say it:} \emph{āsthā}, once more
  \item \emph{Say it:} \emph{merī dādī kī pūjā meṁ gahrī āsthā hai.} — My grandmother has deep faith in worship.
  \item \emph{Recall:} say the Hindi for integrity, then the Hindi for honour, then say \emph{āsthā} again
\end{itemize}

\begin{tcolorbox}[breakable,colback=teal!4,colframe=teal!35!black,title={Before you move on}]
What does \hi{आस्था} mean? (\textquotedblleft{}Belief, steady faith (in something)\textquotedblright{}.) Say it once more.
\end{tcolorbox}
\section[\texorpdfstring{\hi{अपनाना} (apnānā)}{अपनाना (apnānā)}]{\hi{अपनाना} (apnānā) — to adopt, to take up as one's own}

\label{lesson:HI-C357-apnana}



\begin{quote}
Before the new one: say the Hindi for honour, then the Hindi for belief.
\end{quote}

\subsection*{You'll want to know: \hi{अपनाना}}
\textbf{\hi{अपनाना}} — \emph{apnānā} — \textquotedblleft{}to adopt, to take up as one's own\textquotedblright{}.

\begin{quote}
\hi{मीरा} \hi{ने} \hi{गाँव} \hi{की} \hi{भाषा} \hi{अपना} \hi{ली।}
\end{quote}

\begin{quote}
\emph{Mīrā ne gā̃v kī bhāṣā apnā lī.}
\end{quote}

\begin{quote}
Mira took the village's language as her own.
\end{quote}

Past, with \hi{ने} (he or she …): \hi{उसने} \hi{यह} \hi{आदत} \hi{अपनाई}, \emph{usne yah ādat apnāī}.

\begin{cousinweb}
\textbf{\hi{अपनाना}} (\emph{apnānā}) is a verb made on \hi{अपना}, \textquotedblleft{}one's own\textquotedblright{}, which CDIAL traces to a reconstructed \emph{*ātmanaka}, \textquotedblleft{}own\textquotedblright{} (CDIAL 1135): to make something one's own.
\end{cousinweb}

\begin{grammarlens}[title={What you've built}]
One more word for giving reasons and opinions.
\end{grammarlens}

\subsection*{Your turn}
\begin{itemize}
\raggedright
  \item \emph{Say it:} \emph{apnānā}
  \item \emph{Say it:} \emph{apnānā}, once more
  \item \emph{Say it:} \emph{Mīrā ne gā̃v kī bhāṣā apnā lī.} — Mira took the village's language as her own.
  \item \emph{Recall:} say the Hindi for honour, then the Hindi for belief, then say \emph{apnānā} again
\end{itemize}

\begin{tcolorbox}[breakable,colback=teal!4,colframe=teal!35!black,title={Before you move on}]
What does \hi{अपनाना} mean? (\textquotedblleft{}To adopt, to take up as one's own\textquotedblright{}.) Say it once more.
\end{tcolorbox}
